\subsubsection{Flexão}

A tensão de flexão é dada pela equação:

\begin{equation}
  \sigma_f = \frac{M_{f_t}}{W_{f_t}}
  \label{eq:tensao_flexao_tambor}
\end{equation}

Onde:

\begin{itemize}
  \item $\sigma_f$: tensão de flexão no tambor em \si{\mega\pascal}.
  \item $M_{f_t}$: momento fletor máximo no tambor em \si{\newton\meter} (Tabela \ref{tab:dados_esforcos_tambor}).
  \item $W_{f_t}$: módulo de resistência à flexão do tambor em \si{\meter\cubed}.
\end{itemize}

\vspace{1em}

O módulo de resistência à flexão do tambor é dado pela equação:

\begin{equation}
  W_{f_t} = \frac{\pi \left( {D_u}^4 - {D_i}^4 \right)}{32 D_u}
  \label{eq:modulo_resistencia_flexao_tambor}
\end{equation}

Onde:

\begin{itemize}
  \item $W_{f_t}$: módulo de resistência à flexão do tambor em \si{\meter\cubed}.
  \item $D_u$: diâmetro usinado do tambor em \si{\meter} \eqref{eq:diametro_usinado}.
  \item $D_i$: diâmetro interno do tambor em \si{\meter} (Tabela \ref{tab:dados_do_tambor_adotado}).
\end{itemize}

\vspace{1em}

Substituindo os valores na equação \eqref{eq:modulo_resistencia_flexao_tambor}:

\begin{align*}
  W_{f_t} & = \frac{\pi \left( 0,34294^4 - 0,30794^4 \right)}{32 \cdot 0,34294} \\
  W_{f_t} & = \boxed{\SI{1,39e-3}{\meter\cubed}}
\end{align*}

Substituindo os valores na equação \eqref{eq:tensao_flexao_tambor}:

\begin{align*}
  \sigma_f & = \frac{27681,86}{1,39 \cdot 10^{-3}} \\
  \sigma_f & = \boxed{\SI{19,92}{\mega\pascal}}
\end{align*}

\subsubsection{Torção}

A tensão de cisalhamento por torção é dada pela equação:

\begin{equation}
  \tau_T = \frac{T_s}{W_{T_t}}
  \label{eq:tensao_cisalhamento_torcao_tambor}
\end{equation}

Onde:

\begin{itemize}
  \item $\tau_T$: tensão de cisalhamento por torção no tambor em \si{\mega\pascal}.
  \item $T_s$: torque de saída do redutor em \si{\newton\meter} \eqref{eq:torque_saida_redutor}.
  \item $W_{T_t}$: módulo de resistência à torção do tambor em \si{\meter\cubed}.
\end{itemize}

\vspace{1em}

O módulo de resistência à torção do tambor é dado pela equação:

\begin{equation}
  W_{T_t} = \frac{\pi \left( {D_u}^4 - {D_i}^4 \right)}{16 D_u}
  \label{eq:modulo_resistencia_torcao_tambor}
\end{equation}

Onde:

\begin{itemize}
  \item $W_{T_t}$: módulo de resistência à torção do tambor em \si{\meter\cubed}.
  \item $D_u$: diâmetro usinado do tambor em \si{\meter} \eqref{eq:diametro_usinado}.
  \item $D_i$: diâmetro interno do tambor em \si{\meter} (Tabela \ref{tab:dados_do_tambor_adotado}).
\end{itemize}

\vspace{1em}

Substituindo os valores na equação \eqref{eq:modulo_resistencia_torcao_tambor}:

\begin{align*}
  W_{T_t} & = \frac{\pi \left( 0,34294^4 - 0,30794^4 \right)}{16 \cdot 0,34294} \\
  W_{T_t} & = \boxed{\SI{2,77e-3}{\meter\cubed}}
\end{align*}

Substituindo os valores na equação \eqref{eq:tensao_cisalhamento_torcao_tambor}:

\begin{align*}
  \tau_T & = \frac{9150,44}{2,77 \cdot 10^{-3}} \\
  \tau_T & = \boxed{\SI{3,3}{\mega\pascal}}
\end{align*}

\subsubsection{Cisalhamento}

A tensão de cisalhamento cortante é dada pela equação:

\begin{equation}
  \tau_V = \frac{4}{3} \frac{V_t}{A_m}
  \label{eq:tensao_cisalhamento_tambor}
\end{equation}

Onde:

\begin{itemize}
  \item $\tau_V$: tensão de cisalhamento cortante no tambor em \si{\mega\pascal}.
  \item $V_t$: força cortante máxima no tambor em \si{\newton}.
  \item $A_m$: área da seção transversal do tambor em \si{\meter\squared}.
\end{itemize}

\vspace{1em}

A área da seção transversal do tambor é dada pela equação:

\begin{equation}
  A_m = \frac{\pi \left( {D_u}^2 - {D_i}^2 \right)}{4}
  \label{eq:area_secao_transversal_tambor}
\end{equation}

Onde:

\begin{itemize}
  \item $A_m$: área da seção transversal do tambor em \si{\meter\squared}.
  \item $D_u$: diâmetro usinado do tambor em \si{\meter} \eqref{eq:diametro_usinado}.
  \item $D_i$: diâmetro interno do tambor em \si{\meter} (Tabela \ref{tab:dados_do_tambor_adotado}).
\end{itemize}

\vspace{1em}

Substituindo os valores na equação \eqref{eq:area_secao_transversal_tambor}:

\begin{align*}
  A_m & = \frac{\pi \left( 0,34294^2 - 0,30794^2 \right)}{4} \\
  A_m & = \boxed{\SI{17,89e-3}{\meter\squared}}
\end{align*}

Substituindo os valores na equação \eqref{eq:tensao_cisalhamento_tambor}:

\begin{align*}
  \tau_V & = \frac{4}{3} \frac{30320,2}{17,89 \cdot 10^{-3}} \\
  \tau_V & = \boxed{\SI{2,26}{\mega\pascal}}
\end{align*}

\subsubsection{Tensão combinada}

A tensão combinada é dada pela equação:

\begin{equation}
  \sigma_c = \sqrt{{\sigma_f}^2 + 3 {\left( \tau_T + \tau_V \right)}^2}
  \label{eq:tensao_combinada_tambor}
\end{equation}

Onde:

\begin{itemize}
  \item $\sigma_c$: tensão combinada no tambor em \si{\mega\pascal}.
  \item $\sigma_f$: tensão de flexão no tambor em \si{\mega\pascal}.
  \item $\tau_T$: tensão de cisalhamento por torção no tambor em \si{\mega\pascal}.
  \item $\tau_V$: tensão de cisalhamento cortante no tambor em \si{\mega\pascal}.
\end{itemize}

Substituindo os valores na equação \eqref{eq:tensao_combinada_tambor}:

\begin{align*}
  \sigma_c & = \sqrt{{\left( 19,92 \right)}^2 + 3 {\left( 3,3 + 2,26 \right)}^2} \\
  \sigma_c & = \boxed{\SI{22,13}{\mega\pascal}}
\end{align*}

Verificando a tensão combinada com a tensão admissível \eqref{eq:tensao_admissivel_tambor}:

\begin{align*}
  \sigma_c                 & \le \sigma_{\text{adm}}                             \\
  \SI{22,13}{\mega\pascal} & \le \SI{148,21}{\mega\pascal} \rightarrow \text{OK}
\end{align*}
